\documentclass[11pt, a4paper]{../problems_template}

\pagestyle{empty}

\begin{document}

\begin{titlepage}
  \begin{center}
	\Huge{Всероссийская олимпиада}
  \end{center}
\end{titlepage} 

\chapter{1992-1993}

\section{Заключительный этап}

\begin{exercise}
Натуральное число $n$ таково, что числа $2n + 1$ и $3n + 1$ являются
квадратами. Может ли при этом число $5n + 3$ быть простым?
\end{exercise}

\begin{Solution}
Пусть $2n+1=x^{2}$, а $3n+1 = y^{2}$. Заметим, что
$$
5n+3 = 4(2n+1) - (3n+1) = 4x^{2} - y^{2} = (2x - y)(2x + y)
$$
Докажем, что $2x - y \neq 1$. Допустим иное. Тогда
$$
\begin{cases}
2x - y = 1 \\
2x + y = 5n + 3
\end{cases} \implies 2y = 5n + 2 \implies (y - 1)^{2} = 3n+1 - (5n + 2) + 1 = -2n < 0
$$
Что противоречит условию.
\end{Solution}

\section{Окружной этап}

\begin{exercise}
Докажите, что уравнение $x^{3} + y^{3} = 4(x^{2} y + xy^{2} + 1)$ не имеет решений в целых числах.
\end{exercise}

\begin{Solution}
Перепишем уравнение в форме
$$
(x + y)^{3} - 7xy(x + y) = 4
$$
Теперь заметим, что левая часть по модулю 7 не может быть равной 4. 
\end{Solution}

\chapter{1993-1994}

\begin{exercise}
В правильном $(6n+1)$-угольнике $K$ вершин покрашено в красный цвет, а остальные - в синий. Докажите, что количество равнобедренных треугольников с одноцветными вершинами не зависит от способа раскраски.
\end{exercise}

\begin{Solution}
Введем обозначение $N = 6n + 1$. Заметим, что среди треугольников нет равносторонних и отрезок между любыми двумя вершинами принадлежит ровно трем равнобедренным треугольникам. Составим соотношения между количествами треугольников и отрезков. Обозначим $T$ - количество треугольников, а $L$ количество отрезков.
\begin{itemize}
\item Количество синих отрезков, в силу того, что каждый такой отрезок входит либо в синий треугольник (который содержит три таких отрезка), либо в треугольник с двумя синими вершинами (содержит один такой отрезок) - $3L_{b} = 3 T_{bbb} + T_{bb}$.
\item Количество красных отрезков вычисляется аналогично - $3L_{r} = 3T_{rrr} + T_{rr}$.
\item Количество красно-синих отрезков, так как их два в треугольнике с двумя синими вершинами и два в треугольнике с двумя красными вершинами - $3L_{rb} = 2T_{bb}+2T_{rr}$
\end{itemize}
Тогда количество одноцветных треугольников
$$
T_{rrr} + T_{bbb} = \frac{(3L_{b} - T_{bb}) + (3L_{r} - T_{rr})}{3} = L_{b} + L_{r} - \frac{L_{rb}}{2} = \binom{N - K}{2} + \binom{K}{2} - \frac{K(N - K)}{2}
$$
\end{Solution}

\section{Окружной этап}

\begin{exercise}
На совместной конференции партий лжецов и правдолюбов в президиум было избрано 32 человека, которых рассадили в четыре ряда по 8 человек. В перерыве каждый член президиума заявил, что среди его соседей есть представители обеих партий. Известно, что лжецы всегда лгут, а правдолюбы всегда говорят правду. При каком наименьшем числе лжецов в президиуме возможна описанная ситуация? (Два члена президиума являются соседями, если один из них сидит слева, справа, спереди или сзади от другого).
\end{exercise}

\chapter{1994-1995}

\begin{exercise}
Докажите, что для любого натурального числа $a_{1} > 1$ существует возрастающая последовательность натуральных чисел $a_{1}, a_{2}, a_{3}, \dots$ такая, что $a_{1}^{2} + a_{2}^{2} + \dots + a_{k}^{2}$ делится на $a_{1} + a_{2} + \dots + a_{k}$ при всех $k > 1$.
\end{exercise}


\section{Окружной этап}

\begin{exercise}
Можно ли расставить по кругу 1995 различных натуральных чисел так, чтобы для любых двух соседних чисел отношение большего из них к меньшему было простым числом?
\end{exercise}

\begin{Solution}
Обозначим эти числа как $a_{1}, a_{2}, a_{3}, \dots, a_{1995}$ и рассмотрим произведение
$$
\frac{a_{2}}{a_{1}} \frac{a_{3}}{a_{2}}\frac{a_{4}}{a_{3}}\cdots \frac{a_{1}}{a_{1995}}
$$
С одной стороны оно равно единице. А с другой стороны
$$
\frac{p_{1} p_{2} \dots p_{n}}{q_{1}q_{2}\dots q_{m}}
$$
Где и $q_{i}$ и $p_{i}$ простые, а также $n + m = 1995$. Но тогда получается что $n$ и $m$ разной четности. А значит эта дробь никак не может быть равной 1. И расставить числа нельзя.
\end{Solution}


\begin{exercise}
Все стороны и диагонали правильного 12-угольника раскрашиваются в 12 цветов (каждый отрезок - одним цветом). Существует ли такая раскраска, что для любых трех цветов найдутся три вершины, попарно соединенные между собой отрезками этих цветов?
\end{exercise}

\begin{Solution}
 
\end{Solution}

\begin{exercise}
Можно ли в таблице $11 \times 11$ расставить натуральные числа от 1 до 121 так, чтобы числа, отличающиеся друг от друга на единицу, располагались в клетках с общей стороной, а все точные квадраты попали в один столбец?
\end{exercise}

\begin{exercise}
Натуральные числа $m$ и $n$ таковы, что
$$
\lcm(m, n) + \gcd(m, n) = m + n 
$$
Докажите, что одно из чисел $m$ или $n$ делится на другое.
\end{exercise}

\begin{Solution}
Сделаем замену $m = p \cdot \gcd(m, n)$ и $n = q \cdot \gcd(m, n)$. 
$$
\lcm(m, n) + \gcd(m, n) = pq \cdot \gcd(m, n) + gcd(m, n) = (p + q) \cdot \gcd(m, n)
$$
Откуда получаем, что
$$
pq + 1 = p + q \implies (p - 1)(q - 1) = 0
$$
Откуда следует утверждение задачи.
\end{Solution}

\chapter{1995-1996}

\begin{exercise}
Можно ли так расставить фишки в клетках доски $8 \times 8$, чтобы в любых двух столбцах количество фишек было одинаковым, а в любых двух строках - различным?
\end{exercise}

\begin{Solution}
Из условия следует, что общее количество фишек с одной стороны равно $36 - n$, где $n \leqslant 8$, а с другой стороны делится на 8. Из этого следует, что $n = 4$. Из этого получаем требуемую расстановку - ставим в первые 4 строки подряд 8, 7, 6, 5 фишек соответственно, а в остальных строках ставим считая с конца строки подряд 3, 2, 1, 0 фишек соответственно.
\end{Solution}

\begin{exercise}
В Думе 1600 депутатов, которые образовали 16000 комитетов по
80 человек в каждом. Докажите, что найдутся два комитета, имеющие не
менее четырех общих членов.
\end{exercise}

\begin{exercise}
Можно ли прямоугольник $5 \times 7$ покрыть уголками из трех клеток
(то есть фигурками, которые получаются из квадрата $2 \times 2$ удалением одной клетки), не выходящими за его пределы, в несколько слоев так, чтобы каждая клетка прямоугольника была покрыта одинаковым числом клеток, принадлежащих уголкам?
\end{exercise}

\begin{exercise}
Может ли число, получаемое выписыванием в строку друг за другом целых чисел от 1 до $n$ ($n > 1$), одинаково читаться слева направо и справа налево?
\end{exercise}

\begin{exercise}
В строку в неизвестном порядке записаны все целые числа от 1 до
100. За один вопрос про любые 50 чисел можно узнать, в каком порядке относительно друг друга записаны эти 50 чисел. За какое наименьшее число вопросов наверняка можно узнать, в каком порядке записаны все 100 чисел?
\end{exercise}


\chapter{1996-1997}

\begin{exercise}
На доске записаны числа $1, 2, 3, \dots, 1000$. Двое по очереди стирают по одному числу. Игра заканчивается, когда на доске остаются два числа. Если их сумма делится на три, то побеждает тот, кто делал первый ход, \\ если нет - то его партнер. Кто из них выиграет при правильной игре?
\end{exercise}

\begin{Solution}
Выигрышная стратегия для второго игрока такая: если первый игрок стирает число $x$, то второй должен стереть $1001 - x$. Докажем это.
В игре будет сделано $499 \cdot 2$ ходов. Первоначальная сумма всех чисел равна $1 \mod 3$. Каждые два хода сумма уменьшается на $1001 =  2 \mod 3$. Значит в итоге сумма будет равна $1 - 2 \cdot 499 = -1 \mod 3$. И второй игрок выигрывает.
\end{Solution}

\chapter{1998-1999}

\section{Заключительный этап}

\begin{exercise}
Сумма цифр в десятичной записи натурального числа $n$ равна 100, а сумма цифр числа $44n$ равна 800. Чему равна сумма цифр числа $3n$?
\end{exercise}

\begin{Solution}

\end{Solution}

\begin{exercise}
В некоторой группе из 12 человек среди каждых 9 найдутся 5 попарно знакомых. Докажите, что в этой группе найдутся 6 попарно знакомых.
\end{exercise}


\section{Окружной этап}

\begin{exercise}
Клетки квадрата $50 \times 50$ раскрашены в четыре цвета. Докажите, что существует клетка, с четырех сторон от которой (то есть сверху, снизу, слева и справа) имеются клетки одного с ней цвета (не обязательно соседние с этой клеткой).
\end{exercise}

\chapter{1999-2000}

\section{Заключительный этап}

\begin{exercise}
Клетки таблицы $100 \times 100$ окрашены в 4 цвета так, что в любой
строке и в любом столбце ровно по 25 клеток каждого цвета. Докажите, что найдутся две строки и два столбца, все четыре клетки на пересечении которых окрашены в разные цвета.
\end{exercise}

\section{Окружной этап}

\begin{exercise}
Среди пяти внешне одинаковых монет 3 настоящие и две фальшивые, одинаковые по весу, но неизвестно, тяжелее или легче настоящих. Как за наименьшее число взвешиваний найти хотя бы одну настоящую монету?
\end{exercise}

\chapter{2000-2001}

\section{Заключительный этап}

\begin{exercise}
Числа от 1 до 999 999 разбиты на две группы: в первую отнесено каждое число, для которого ближайшим к нему квадратом является квадрат нечетного числа, во вторую - числа, для которых ближайшими являются квадраты четных чисел. В какой из групп сумма чисел больше?
\end{exercise}

\section{Окружной этап}

\begin{exercise}
Можно ли клетки доски $5 \times 5$ покрасить в 4 цвета так, чтобы клетки,
стоящие на пересечении любых двух строк и любых двух столбцов, были
покрашены не менее чем в 3 цвета?
\end{exercise}

\begin{Solution}
Рассмотрим, для примера, первую строку. Поскольку в строке 5 клеток, а всего цветов 4, то какой-то цвет появляется минимум два раза. Выберем две столбца которым принадлежат эти клетки. Поскольку для выбранных клеток первой строки мы имеем покраску только в один цвет, то в остальных строках в выбранных столбцах должны быть сочетания двух цветов, отличных от первого. Всего существует $\binom{3}{2} = 3$ способов выбрать два разных цвета из оставшихся трех цветов. Но оставшихся строк четыре. Значит какое то сочетание цветов появляется два раза. Но это означает, что 4 клетки будут окрашены только в 2 цвета.  
\end{Solution}

\begin{exercise}
Докажите, что в любом множестве, состоящем из 117 попарно различных трехзначных чисел, можно выбрать 4 попарно непересекающихся подмножества, суммы чисел в которых равны.
\end{exercise}

\begin{Solution}
Всего существует $\binom{117}{2}$ подмножеств из двух элементов в заданном условием множестве. По условию среди них не может быть более трех с равной суммой. Минимальная сумма двух элементов в множестве равна $100 + 101 = 201$, максимальная равна $998 + 999 = 1997$. Значит сумма может принимать не более $1997 - 201 + 1 = 1797$ значений. Если каждое значение принимается не более трех раз то всего получается $5391 < \binom{117}{2}$.
\end{Solution}

\chapter{2001-2002}

\section{Окружной этап}

\begin{exercise}
Среди 18 деталей, выставленных в ряд, какие-то три подряд стоящие весят по 99 г, а все остальные - по 100 г. Двумя взвешиваниями на весах со стрелкой определите все 99-граммовые детали.
\end{exercise}

\begin{exercise}
В какое наибольшее число цветов можно раскрасить все клетки доски размера $10 \times 10$ так, чтобы в каждой строке и в каждом столбце находились клетки не более, чем пяти различных цветов?
\end{exercise}

\chapter{2002-2003}

\section{Окружной этап}

\begin{exercise}
На встречу выпускников пришло 45 человек. Оказалось, что любые двое из них, имеющие одинаковое число знакомых среди пришедших, не знакомы друг с другом. Какое наибольшее число пар знакомых могло быть среди участвовавших во встрече?
\end{exercise}

\begin{exercise}
В наборе из 17 внешне одинаковых монет две фальшивых, отличающихся от остальных по весу. Известно, что суммарный вес двух фальшивых монет вдвое больше веса настоящей. Всегда ли можно ли определить пару фальшивых монет, совершив 5 взвешиваний на чашечных весах без гирь? (Определить, какая из фальшивых тяжелее, не требуется.)
\end{exercise}


\chapter{2003-2004}

\section{Заключительный этап}

\begin{exercise}
Натуральные числа от 1 до 100 расставлены по кругу в таком порядке, что каждое число либо больше обоих соседей, либо меньше обоих соседей. Пара соседних чисел называется хорошей, если при выкидывании этой пары вышеописанное свойство сохраняется. Какое минимальное
количество хороших пар может быть?
\end{exercise}

\begin{Solution}
\end{Solution}

\section{Окружной этап}

\chapter{2005-2006}

\begin{exercise}
Натуральные числа от 1 до 200 разбили на 50 множеств. Докажите, что в одном из них найдутся три числа, являющиеся длинами сторон некоторого треугольника.
\end{exercise}

\begin{Solution}
Рассмотрим числа от 100 до 200. Их всего 101, а значит три из них попадут в одно множество. Но любые два числа в этом промежутке в сумме больше 200, а значит больше третьего числа. А значит эти три числа являются сторонами треугольника.
\end{Solution}

\begin{exercise}
В лагерь приехало несколько пионеров, каждый из них имеет от 50 до 100 знакомых среди остальных. Докажите, что пионерам можно выдать пилотки, покрашенные в 1331 цвет так, чтобы у знакомых каждого пионера были пилотки хотя бы 20 различных цветов. 
\end{exercise}

\begin{exercise}
Клетчатый квадрат $100 \times 100$ разрезан на доминошки: прямоугольники $1 \times 2$. Двое играют в игру. Каждым ходом игрок склеивает две соседних по стороне клетки, между которыми был проведен разрез. Игрок
проигрывает, если после его хода фигура получилась связной, то есть весь
квадрат можно поднять со стола, держа его за одну клетку. Кто выиграет
при правильной игре - начинающий или его соперник?
\end{exercise}

\begin{exercise}
Квадрат $3000 \times 3000$ произвольным образом разбит на доминошки
(то есть прямоугольники $1 \times 2$ клетки). Докажите, что доминошки можно раскрасить в три цвета так, чтобы доминошек каждого цвета было поровну и у каждой доминошки было не более двух соседей ее цвета (доминошки считаются соседними, если они содержат клетки, соседние по стороне).
\end{exercise}

\chapter{2006-2007. Номер 33}

\begin{exercise}
Среди натуральных чисел от 1 до 1200 выбрали 372 различных числа так, что никакие два из них не различаются на 4, 5 или 9. Докажите, что число 600 является одним из выбранных.
\end{exercise}

\begin{Solution}
Докажем следующую лемму - из 13 подряд идущих чисел можно выбрать не более 4 так, чтобы никакие из них не различались на 4, 5 или 9. Разобьем числа от $a$ до $a + 12$ на группы.
$$
(a + 4), (a, a + 9), (a + 5), (a + 1, a + 10), (a + 6), (a + 2, a + 11), (a + 7), (a + 3, a + 12), (a + 8) 
$$
Заметим, что из каждой группы можно выбрать не более одного элемента и в двух соседних группах в сумме может быть выбрано не более одного элемента. Но если мы можем выбрать хотя бы 5 элементов, то по принципу Дирихле найдутся два элемента в соседних группах. Противоречие, а значит всего элементов не более четырёх. Лемма доказана. Теперь исключим числа 599, 600, 601, 602. Оставшиеся числа составят $\frac{598}{13} + \frac{1200 - 603 + 1}{13} = 92$ группы по 13 чисел.  В них, исходя из леммы, может быть выбрано $92 \cdot 12 = 368$ чисел. Значит числа 599, 600, 601, 602 тоже выбраны. 
\end{Solution}

\begin{exercise}
В таблице $10 \times 10$ расставлены числа от 1 до 100: в первой строчке - от 1 до 10 слева направо, во второй - от 11 до 20 слева направо
и так далее. Андрей собирается разрезать таблицу на прямоугольники $1 \times 2$, посчитать произведение чисел в каждом прямоугольнике и сложить полученные 50 чисел. Он стремится получить как можно меньшую сумму. Как ему следует разрезать квадрат?
\end{exercise}

\chapter{2007-2008. Номер 34}

\begin{exercise}
Существуют ли такие 14 натуральных чисел, что при увеличении каждого из них на 1 произведение всех чисел увеличится ровно в 2008 раз?
\end{exercise}

\begin{exercise}
На доске написано натуральное число. Если на доске написано число $x$, то можно дописать на нее число $2x + 1$ или $\frac{x}{x + 2}$. В какой-то момент выяснилось, что на доске присутствует число 2008. Докажите, что оно там было с самого начала.
\end{exercise}

\section{Окружной этап}

\begin{exercise}
300 бюрократов разбиты на три комиссии по 100 человек. Любые два бюрократа либо знакомы друг с другом, либо незнакомы. Докажите, что найдутся два таких бюрократа из разных комиссий, что в третьей комиссии есть либо 17 человек, знакомых с обоими, либо 17 человек, незнакомых с обоими.
\end{exercise}

\begin{Solution}
Рассмотрим граф знакомств бюрократов находящихся в разных комиссиях. Пусть два знакомых бюрократа будут соединены красным ребром, а два незнакомых - синим. Всего в этом графе $100 \cdot 100 \cdot 100$ треугольников. У каждого треугольника два ребра одного цвета. Значит существует не менее $\frac{100^{3}}{2}$ треугольников с двумя ребрами определенного цвета. Общая вершина этих ребер лежит в одной из трех групп. Значит из какой то группы исходит не менее $\frac{100^{3}}{6}$ пар ребер выбранного цвета. В двух других группах $100^{2}$ пар вершин куда могут идти эти пары ребер. Значит не менее $\frac{100}{6} = 17$ пар ребер приходит к какой-то паре вершин из двух комиссий.
\end{Solution}

\chapter{2008-2009}

\begin{exercise}
По кругу стоят 100 наперстков. Под одним из них спрятана монетка. За один ход разрешается перевернуть четыре наперстка и проверить, лежит ли под одним из них монетка. После этого их возвращают в исходное положение, а монетка перемещается под один из соседних с ней наперстков. За какое наименьшее число ходов наверняка удастся обнаружить монетку?
\end{exercise}

\begin{exercise}
Можно ли раскрасить натуральные числа в 2009 цветов так, чтобы каждый цвет встречался бесконечное число раз, и не нашлось тройки чисел, покрашенных в три различных цвета, таких, что произведение двух из них равно третьему?
\end{exercise}

\begin{Solution}
Возможно несколько вариантов такой раскраски.
\begin{enumerate}
\item Раскрасим все простые числа так, что $p_{n}$ окрашивается в цвет $n \mod 2009$. Цвет числа $N = p_{1}^{\alpha_{1}} p_{2}^{\alpha_{2}} \dots p_{k}^{\alpha_{k}}$, где $p_{1} < p_{2} < \dots < p_{k}$ тот же, что у $p_{1}$. Тогда если два числа, $N$ и $M$ имеют разные цвета, то их произведение будет иметь либо цвет $N$ либо цвет $M$.
\end{enumerate}
\end{Solution}

\begin{exercise}
Последовательность $a_{1}, a_{2}, \dots$ такова, что $a_{1} \in (1, 2)$ и $a_{k+1} = a_{k} + \frac{k}{a_{k}}$ при любом натуральном $k$. Докажите, что в ней не может существовать более одной пары членов с целой суммой.
\end{exercise}

\chapter{2013-2014}

\begin{exercise} 
Назовем натуральное число хорошим, если среди его делителей есть ровно два простых числа. Могут ли 18 подряд идущих натуральных чисел быть хорошими?
\end{exercise}

\chapter{2014-2015}

\begin{exercise}
Назовём натуральное число почти квадратом, если оно равно произведению двух последовательных натуральных чисел. Докажите, что каждый почти квадрат можно представить в виде частного двух почти квадратов.
\end{exercise}


\chapter{2019-2020}

\begin{exercise}
Множество $A$ состоит из $n$ различных натуральных чисел, сумма которых равна $n^{2}$. Множество $B$ также состоит из $n$ различных натуральных чисел, сумма которых равна $n^{2}$. Докажите, что найдётся число, которое принадлежит как множеству $A$, так и множеству $B$.
\end{exercise}

\begin{Solution}
Предположим противное. Рассмотрим сумму всех чисел в обоих множествах. Поскольку все числа разные, то она не менее $n \cdot (2n + 1)$. Но по условию она равна $2n^{2}$. Противоречие, а значит множества пересекаются.
\end{Solution}

\chapter{2020-2021. Номер 47}

\section{Заключительный этап}

\begin{exercise}
На окружности отмечено 1000 точек, каждая окрашена в один из $k$ цветов. Оказалось, что среди любых пяти попарно пересекающихся отрезков, концами которых являются 10 различных отмеченных точек, найдутся хотя бы три отрезка, у каждого из которых концы имеют разные цвета. При каком наименьшем $k$ это возможно?
\end{exercise}

\begin{Solution}
Допустим существует 8 точек одного цвета, Обозначим их $A_{1}, A_{2}, A_{3}, \dots, A_{8}$ по часовой стрелке. Тогда взяв отрезки $A_{1}A_{6}$, $A_{2}A_{7}$, $A_{3}A_{8}$ мы получим три отрезка одного цвета. И отрезков с концами в точках $A_{4}$ и $A_{5}$, которые вместе с существующими сформируют конструкцию из условия задачи, не хватит, чтобы  среди попарно пересекающихся 5 отрезков было три отрезка, у каждого из которых концы имеют разные цвета. Значит точек каждого цвета не более 7. Значит $k \geqslant \frac{1000}{7}$, то есть $k \geqslant 143$.
\end{Solution}

\chapter{2022-2023}

\section{Заключительный этап}

\begin{exercise}
Каждое натуральное число, большее 1000, окрасили либо в красный, либо в синий цвет. Оказалось, что \\ произведение любых двух различных красных чисел - синее. Может ли случиться, что никакие два синих числа не отличаются на 1?
\end{exercise}

\begin{exercise}
У 100 школьников есть стопка из 101 карточки, которые пронумерованы числами от 0 до 100. Первый школьник перемешивает стопку, затем берёт сверху из получившейся стопки по одной карточке, и при каждом взятии карточки (в том числе при первом) записывает на доску среднее арифметическое чисел на всех взятых им на данный момент карточках. Так он записывает 100 чисел, а когда в стопке остаётся одна карточка, он возвращает карточки в стопку, и далее всё то же самое, начиная с перемешивания стопки,
проделывает второй школьник, потом третий, и так далее. Докажите, что среди выписанных на доске 10000 чисел найдутся два одинаковых.
\end{exercise}

\begin{exercise}
Рассмотрим все 100-значные числа, делящиеся на 19. Докажите, что количество таких чисел, не содержащих цифр 4, 5 и 6, равно количеству таких чисел, не содержащих цифр 1, 4 и 7.
\end{exercise}

\section{Региональный этап}

\begin{exercise}
В таблице $6 \times 6$ изначально записаны нули. За одну операцию можно выбрать одну клетку и заменить число, стоящее в ней, на любое целое число. Можно ли за 8 операций получить таблицу,
в которой все 12 сумм чисел в строках и столбцах будут различными положительными числами?
\end{exercise}

\begin{exercise}
В городе $N$ прошли 50 городских олимпиад по разным предметам, при этом в каждой из этих олимпиад участвовало ровно 30 школьников, но не было двух олимпиад с одним и тем же составом участников. Известно, что для любых 30 олимпиад найдется школьник, который участвовал во всех этих 30 олимпиадах. Докажите, что найдется школьник, который участвовал во всех 50 олимпиадах.
\end{exercise}


\chapter{2023-2024}

\begin{exercise}

\end{exercise}



\chapter{2024-2025}

\begin{exercise}
На прямоугольном листе бумаги провели несколько отрезков, параллельных его сторонам. Эти отрезки разбили лист на несколько прямоугольников, внутри которых нет проведённых линий. Петя хочет провести в каждом из прямоугольников разбиения одну диагональ, разбив его на два треугольника, и окрасить каждый треугольник либо в чёрный, либо в белый цвет. Верно ли, что он обязательно сможет это сделать так, чтобы никакие два одноцветных треугольника не имели общего отрезка границы?
\end{exercise}

\begin{Solution}
\end{Solution}

\begin{exercise}
Найдите все натуральные $n$, для которых существует такое чётное натуральное $a$, что число 
$$
(a - 1)(a^{2} - 1)\dots(a^{n} - 1)
$$ 
является точным квадратом.
\end{exercise}

\begin{exercise}
Шахматного короля поставили на клетку доски $8 \times 8$ и сделали им 64 хода так, что он побывал на всех клетках и вернулся в исходную клетку. В каждый момент времени вычислялось расстояние от центра клетки, в которой находился король, до центра всей доски. Назовём сделанный ход приятным, если в результате хода это расстояние стало меньше, чем было до хода. Найдите наибольшее возможное
количество приятных ходов. (Шахматный король за один ход передвигается на клетку, соседнюю по стороне или по углу.)
\end{exercise}

\begin{exercise}
Петя и Вася играют в игру на изначально пустой клетчатой таблице $100 \times 100$, делая ходы по очереди. Начинает Петя. За свой ход игрок вписывает в некоторую пустую клетку любую заглавную букву русского алфавита (в каждую клетку можно вписать ровно одну букву). Когда все клетки будут заполнены, Петя объявляется победителем, если найдутся четыре подряд идущие клетки по горизонтали, в
которых слева направо написано слово "ПЕТЯ", или найдутся четыре подряд идущие клетки по вертикали, в которых сверху вниз написано слово "ПЕТЯ". Сможет ли Петя
выиграть независимо от действий Васи?
\end{exercise}

\begin{exercise}
На доску выписали 777 попарно различных комплексных чисел. Оказалось, что можно ровно 760 способами выбрать два числа $a$ и $b$, записанных на доске, так, чтобы выполнялось равенство 
$$a^{2} + b^{2} + 1 = 2ab
$$
Способы, которые отличаются перестановкой чисел, считаются одинаковыми. Докажите, что можно выбрать такие
два числа $c$ и $d$, записанных на доске, что
$$
c^{2} + d^{2} + 2025 = 2cd
$$
\end{exercise}

\begin{Solution}
Можно заметить, что первое равенство переписывается в виде
$$
(a - b)^{2} = -1 \implies |a - b| = i
$$
Обозначим количество вершин в цепи $i$ как $v_{i}$, а количество ребер как $e_{i}$. Тогда
$$
\sum_{i = 1}^{N} v_{i} = 777, \quad \sum_{i = 1}^{N} e_{i} = 760, \quad v_{i} - e_{i} = 1
$$
\end{Solution}

\begin{exercise}
Петя выбрал 100 попарно различных положительных чисел, меньших 1, и расставил их по кругу. Затем он проделывает с ними операции. За одну операцию можно взять
три стоящих подряд (именно в таком порядке) числа $a$, $b$, $c$ и заменить число $b$ на $a - b + c$. При каком наибольшем $k$ Петя мог выбрать исходные числа и сделать несколько операций так, чтобы после них среди чисел оказалось $k$ целых?
\end{exercise}


\end{document} 
 
 
 
